%!TEX root = thesis.tex
\section{Artificial neurons}\label{sec:artificial-neurons}
%% ===========================

Artificial neurons are inspired by biological neurons.

\begin{figure}[ht]
    \centering
    \subfloat[Biological Neuron]{
        \resizebox{0.45\linewidth}{!}{\begin{tikzpicture}
    \node[draw,circle] (x) at (0,0) {TODO};
\end{tikzpicture}}
        \label{fig:biological-neuron}
    }%
    \subfloat[Artificial Neuron]{
        \resizebox{0.45\linewidth}{!}{\input{figures/artificial-neuron.tex}}
        \label{fig:artificial-neuron}
    }%
    \label{fig:artificial-and-biological-neuron}
    \caption{Both neurons receive weighted input, apply a function to that and give output}
\end{figure}

They receive at least one
input and at least one output. Those inputs might get weighted as well as the
output.

The neurons apply a function to the sum of all weighted inputs. This function
is called \textit{activaiton function}.

An artificial neuron using the unit step function (see \cref{f:unitstep}) is called
a \textit{perceptron}.


The artificial neuron sums all weighted inputs $x_i \cdot w_i$ up
             and applies its activation function $f$ to it.